\newcommand\Fname{Роман}
\newcommand\Lname{Уколов}
\newcommand\Mname{Олегович}
\newcommand\Group{ІМ-41}
\newcommand\Nlist{22}

\newcommand\Department{обчислювальної техніки}
\newcommand\Variant{22}

\usepackage{xstring}

\newtoggle{male}
\newtoggle{female}

% ця команда отримує позначення дисципліни та обирає приймача роботи.
% я написав сценарій, який визначає це позначення за шляхом до поточного каталогу
% й після ще кількох перевірок генерує шаблон звіту, де заповнене все, крім дати й теми.
% Звісно, не завжди такий шаблонний метод працює, але тоді можна використати \renewcommand
\newcommand{\Work}[1]{
\IfStrEqCase{#1}{
    {os}{
        \newcommand\discipline{Операційні системи}
        \IfStrEqCase{\Type}{
            {lab}{
                \newcommand\instructor{Андрій Сімоненко}
                \toggletrue{male}
            }
        }
    }
    {linux}{
        \newcommand\discipline{Linux}
        \IfStrEqCase{\Type}{
            {lab}{
                \newcommand\instructor{Марина Хмелюк}
                \toggletrue{female}
                \newcommand\academicRank{ст. вик. кафедри ІСТ}
            }
        }
    }
    {<++>}{
        \newcommand\discipline{<++>}
        \IfStrEqCase{\Type}{
            {lab}{
                \newcommand\instructor{<++>}
            }
        }
    }
}
}
